\documentclass[11pt,a4paper]{article}
\usepackage[T1]{fontenc}
\usepackage{lmodern}
\usepackage[margin=25mm]{geometry}
\usepackage{amsmath,amssymb,amsthm,booktabs,microtype,enumitem}
\usepackage[hidelinks]{hyperref}
\hypersetup{pdftitle={Congruences for the scaled inverse series A393859},pdfauthor={Dev Goyal},pdfsubject={A393859; formal power series and Lean proofs}}
\setlength{\parskip}{3pt}
\setlength{\emergencystretch}{3em}
\newtheorem{theorem}{Theorem}[section]
\newtheorem{lemma}[theorem]{Lemma}
\newtheorem{corollary}[theorem]{Corollary}
\theoremstyle{remark}\newtheorem{remark}[theorem]{Remark}
\newcommand{\Z}{\mathbb Z}
\newcommand{\N}{\mathbb N}
\newcommand{\F}{\mathbb F}
\newcommand{\coeff}[2]{[x^{#1}]#2}
\title{Congruences for the scaled inverse series A393859}
\author{Dev Goyal}
\date{24 October 2026}
\begin{document}
\maketitle
\begin{center}\fbox{\begin{minipage}{0.91\linewidth}
\textbf{AI disclosure and review status.}
ChatGPT (OpenAI) generated the mathematical arguments, formal proof code,
and this manuscript from Dev Goyal's prompts. Dev Goyal is the sole human
author and initiated the investigation. This consolidated presentation is
AI-generated; independent human review of this presentation is pending.
Source conjectures and prior results retain the credits given below.
\end{minipage}}\end{center}

\begin{abstract}
For the scaled inverse series A393859 we prove that every
positive-index coefficient is one modulo 8.
The source equation, coefficientwise interpretation of its quotient,
existence, uniqueness, and the all-index proof are given explicitly.
\end{abstract}
\section{Source equation and theorem}
For a positive integer $q$, consider the normalized series $G_q$ satisfying
\begin{equation}\label{eq:scaled-source}
 G_q\left(x-\frac{xG_q(qx)}q\right)=x,
 \qquad G_q=x+O(x^2).
\end{equation}
The source A393859 of Paul D. Hanna is the case $q=7$
\cite{oeis}. Put $a_n=[x^n]G_{7}$. Its source indices
start at one and $a_0=0$. We prove
\[
 a_n\equiv1\pmod{8}\qquad(n\ge1).
\]

All series are formal: equality and congruence are coefficientwise, and
composition $U(V)=\sum_{j\ge0}u_jV^j$ is used only when $V(0)=0$.
Every coefficient then receives finitely many contributions.
A series with constant coefficient one is a multiplicative unit:
if $W(0)=0$, its inverse is $\sum_{j\ge0}(-W)^j$.
This justifies the rational comparison series below without analytic
convergence or division in a residue ring.
\subsection{Removing the apparent division}

If $U(x)=\sum_{n\ge1}u_nx^n$, define
\[
L_q(U)(x)=\sum_{n\ge1}q^{n-1}u_nx^n.
\]
For a nonzero integer $q$, this is exactly $U(qx)/q$. The coefficient
definition makes sense over every commutative ring, without requiring
that $q$ be invertible. Equation~\eqref{eq:scaled-source} therefore
means
\begin{equation}\label{eq:scaled-nodiv}
G_q\bigl(x-xL_q(G_q)(x)\bigr)=x.
\end{equation}

\begin{lemma}\label{lem:scaled-unique}
For any commutative ring $R$ with identity and any $q\in R$, equation
\eqref{eq:scaled-nodiv} has a unique normalized solution in $R[[x]]$.
Its construction commutes with ring homomorphisms.
\end{lemma}

\begin{proof}
Write $G=x+\sum_{n\ge2}g_nx^n$ and
$V=x-xL_q(G)$. The series $V$ is normalized. The coefficient $g_n$
inside $V$ first occurs in degree $n+1$, because of the initial
factor $x$. Thus the coefficient of degree $n$ in $V$ depends
only on $g_2,\ldots,g_{n-1}$.

In $G(V)=V+\sum_{j\ge2}g_jV^j$, the term $g_nV^n$ contributes
exactly $g_n$ in degree $n$. The terms with $j<n$ depend only on
earlier coefficients, and those with $j>n$ contribute nothing in
that degree. Therefore the equation in degree $n$ has the form
$g_n+P_n=0$, with $P_n$ determined by earlier coefficients and $q$.
This gives existence and uniqueness by recursion, without division.
It also shows that applying a ring homomorphism to the coefficients
preserves the equation; uniqueness identifies the resulting solution.
\end{proof}

\subsection{The congruence and parity theorems}

\begin{theorem}\label{thm:scaled-congruence}
For every positive integer $q$,
\[
G_q(x)\equiv\frac{x}{1-x}\pmod{q+1}.
\]
In particular, every coefficient of positive degree is congruent to
one modulo $q+1$.
\end{theorem}

\begin{proof}
Work in $R=\Z/(q+1)\Z$, where $q=-1$. For $T=x/(1-x)$,
\[
L_{-1}(T)=x-x^2+x^3-\cdots=\frac{x}{1+x}.
\]
Hence
\[
x-xL_{-1}(T)=x-\frac{x^2}{1+x}=\frac{x}{1+x},
\]
and direct substitution gives
\[
T\left(\frac{x}{1+x}\right)
=\frac{x/(1+x)}{1-x/(1+x)}=x.
\]
Thus $T$ is the normalized solution of the reduced equation.
Lemma~\ref{lem:scaled-unique} identifies it with the reduction of
$G_q$. This argument also covers composite moduli such as six
and eight.
\end{proof}

\begin{theorem}\label{thm:scaled-parity}
If $q$ is a positive even integer, then
\[
G_q(x)\equiv\sum_{j\ge0}x^{2^j}\pmod2.
\]
Equivalently, its coefficient of degree $n\ge0$ is odd if and only
if $n$ is a power of two. If $q$ is odd, all positive-index
coefficients of $G_q$ are odd.
\end{theorem}

\begin{proof}
For even $q$, reduction modulo two gives $L_q(U)=x$ for every
normalized $U$: all factors $q^{n-1}$ with $n\ge2$ vanish.
Therefore the reduced equation is
\[
G_q(x+x^2)=x.
\]
Let $S(x)=\sum_{j\ge0}x^{2^j}$ over $\Z/2\Z$.
The strictly increasing exponents make this a well-defined formal
series. Repeated squaring in characteristic two gives
$(x+x^2)^{2^j}=x^{2^j}+x^{2^{j+1}}$, and hence
\[
S(x+x^2)
=\sum_{j\ge0}\bigl(x^{2^j}+x^{2^{j+1}}\bigr)=x.
\]
Every term beyond the first occurs twice and cancels. By uniqueness,
the reduction of $G_q$ is $S$. Its support is exactly the powers
of two, including $1=2^0$ and excluding zero.

If $q$ is odd, then $q+1$ is even, so
Theorem~\ref{thm:scaled-congruence} already gives oddness of every
positive-index coefficient.
\end{proof}


\section{Application to A393859}
Take $q=7$ in Theorem~\ref{thm:scaled-congruence}.
Its modulus is $q+1=8$, so every positive-index
coefficient has the asserted residue. Since this modulus is even, the same formula also implies
that every positive-index coefficient is odd. The source congruence
is valid at composite modulus 8; the proof used units
with constant coefficient one and did not cancel a nonunit parameter.

\section{Formal proof and scope}
The accompanying entry file is \path{lean/Entries/A393859.lean}.
It states the results with explicit Lean types and supplies proof terms
using the complete shared implementation included in the collection.
The manuscript is self-contained mathematically; the Lean entry file uses
those ordinary library imports rather than duplicating the full foundation.

The entry's named theorem interfaces are listed below. The namespace of
each is \path{FinalCollection.A393859}.
\begin{itemize}[leftmargin=1.5em,itemsep=1pt]
\item \path{source_existence_uniqueness}
\item \path{source_normalized}
\item \path{source_equation}
\item \path{a393859_mod_eight}
\end{itemize}
The preserved dependencies also contain all supporting constructions,
coefficient identities, map-compatibility results, and inverse proofs.
The collection pins Lean 4.19.0 and Mathlib revision
\path{c44e0c8ee63ca166450922a373c7409c5d26b00b}.
The verification report records fresh compilation and the transitive axiom
audit. Finite numerical checks corroborate the definitions and indexing;
the universally quantified Lean theorems establish the all-index claims.
No assertion of worldwide novelty, OEIS acceptance, publication, or completed
human review follows from these checks. The prose may use elementary
identities that the formal proof establishes by a different algebraic route.

\begin{thebibliography}{9}\raggedright
\bibitem{oeis} P. D. Hanna, \emph{A393859}, The On-Line Encyclopedia of Integer
Sequences. \url{https://oeis.org/A393859}. Source recorded in the retained
proof collection, 24 October 2026.
\bibitem{repo} D. Goyal, \emph{OEIS A396843 proof repository}.
\url{https://github.com/dvgyl/oeis-a396843-proof}.
Reference revision \path{995e55adcbe06b2a913e24f12ff73a7faf362f06}.
\end{thebibliography}
\end{document}
